\documentclass{article}
\usepackage{amsmath}
\usepackage{amssymb}
\usepackage{amsthm}
\usepackage{mathtools}
\usepackage{geometry}
\usepackage{hyperref}

\newcommand{\rk}{\operatorname{rank}}
\newcommand{\GL}{\operatorname{GL}}
\newcommand{\R}{\mathbb{R}}
\newcommand{\Mat}{\mathrm{Mat}}
\newcommand{\Rep}{\operatorname{Rep}}

\theoremstyle{definition}
\newtheorem{theorem}{Theorem}[section]
\newtheorem{proposition}[theorem]{Proposition}
\newtheorem{lemma}[theorem]{Lemma}
\newtheorem{corollary}[theorem]{Corollary}
\newtheorem{definition}[theorem]{Definition}
\newtheorem{remark}[theorem]{Remark}

\title{Task 4: Krull--Schmidt and Multiplicity Patterns}
\date{}

\begin{document}
\maketitle

\section{Goal}

This file passes from the indecomposable interval modules to arbitrary modules
over
\[
A_L^{\mathrm{wt}}.
\]

The main deliverables are:

\begin{itemize}
\item the Krull--Schmidt decomposition in terms of the admissible interval
  modules $M(i,j)$;
\item the exact multiplicity-pattern set
  $\mathcal M_d^+(A_L^{\mathrm{wt}})$;
\item the orbit dictionary that will later label the cyclic locus.
\end{itemize}

\section{Krull--Schmidt}

\begin{proposition}[Krull--Schmidt for $A_L^{\mathrm{wt}}$]
Every finite-dimensional $A_L^{\mathrm{wt}}$-module decomposes as a finite
direct sum of indecomposable modules, and this decomposition is unique up to
reordering and isomorphism of the summands.
\end{proposition}

\begin{proof}
The algebra $A_L^{\mathrm{wt}}$ is finite-dimensional, hence artinian.
Therefore the category of finite-dimensional modules over
$A_L^{\mathrm{wt}}$ is a Krull--Schmidt category.
\end{proof}

\begin{corollary}[Decomposition in interval modules]
Every finite-dimensional $A_L^{\mathrm{wt}}$-module $M$ has a unique
decomposition
\[
M
\cong
\bigoplus_{0\le i\le j\le L} M(i,j)^{\oplus m(i,j)}
\oplus
\bigoplus_{\substack{1\le i\le L\\0\le j\le i-2}}
M(i,j)^{\oplus m(i,j)},
\]
for uniquely determined multiplicities
\[
m(i,j)\in \mathbb Z_{\ge 0}.
\]
\end{corollary}

\begin{remark}
The first sum is the non-wrapping type-$A$ contribution.
The second sum is the genuinely cyclic contribution coming from the closing
arrow.
\end{remark}

\section{Multiplicity patterns}

\begin{definition}[Support indicator]
For an admissible pair $(i,j)\in \mathcal I_L^{\mathrm{wt}}$ and a vertex
$k\in\{0,\dots,L\}$, define
\[
\mathbf 1_{[i\to j]}(k)
:=
\begin{cases}
1, & k\in [i\to j],\\
0, & k\notin [i\to j].
\end{cases}
\]
\end{definition}

\begin{definition}[Multiplicity-pattern set]
Fix a dimension vector
\[
d=(d_0,\dots,d_L).
\]
Define
\[
\mathcal M_d^+(A_L^{\mathrm{wt}})
:=
\left\{
(m(i,j))_{(i,j)\in \mathcal I_L^{\mathrm{wt}}}
\in
\mathbb Z_{\ge 0}^{\mathcal I_L^{\mathrm{wt}}}
\ \middle|\
d_k=
\sum_{(i,j)\in \mathcal I_L^{\mathrm{wt}}}
m(i,j)\,\mathbf 1_{[i\to j]}(k)
\text{ for all }k
\right\}.
\]
\end{definition}

\begin{proposition}[Exact reconstruction of the dimension vector]
Let
\[
M
\cong
\bigoplus_{(i,j)\in \mathcal I_L^{\mathrm{wt}}}
M(i,j)^{\oplus m(i,j)}.
\]
Then the dimension vector of $M$ is given by
\[
d_k=
\sum_{(i,j)\in \mathcal I_L^{\mathrm{wt}}}
m(i,j)\,\mathbf 1_{[i\to j]}(k),
\qquad
0\le k\le L.
\]
\end{proposition}

\begin{proof}
Each indecomposable interval module contributes a one-dimensional space at
exactly the vertices in its support interval and contributes zero elsewhere.
Summing these contributions over the direct-sum decomposition gives the stated
formula.
\end{proof}

\section{Orbit dictionary}

\begin{proposition}[Orbits and isomorphism classes]
Fix the representation space
\[
\Rep_d(A_L^{\mathrm{wt}})
\]
and the base-change group
\[
G_d:=\prod_{\ell=0}^L \GL(d_\ell).
\]
Then the $G_d$-orbits in $\Rep_d(A_L^{\mathrm{wt}})$ are exactly the
isomorphism classes of $A_L^{\mathrm{wt}}$-modules of dimension vector $d$.
\end{proposition}

\begin{corollary}[Multiplicity patterns parametrize the orbits]
There is a canonical bijection
\[
G_d\backslash \Rep_d(A_L^{\mathrm{wt}})
\cong
\mathcal M_d^+(A_L^{\mathrm{wt}}),
\]
obtained by sending a module to the multiplicities of its indecomposable
interval summands.
\end{corollary}

\begin{remark}
This is the combinatorial data we will later transport from the cyclic locus
to the critical locus by using the orbit correspondence of the first
conjecture.
\end{remark}

\section{Deliverables}

This task is complete once the file contains:

\begin{itemize}
\item the precise Krull--Schmidt decomposition in admissible interval modules;
\item the exact definition of $\mathcal M_d^+(A_L^{\mathrm{wt}})$;
\item the fact that multiplicity patterns parametrize the
  $G_d$-orbits of $\Rep_d(A_L^{\mathrm{wt}})$.
\end{itemize}

\end{document}
